%%% ВЫБОР УСТРОЙСТВА %%%

\IfStrEqCase{\devicetype}{%
{ЮНИТ-М500-ЛВ}{%
% код для этого типа
\def\Qthree{В3}
\def\Qfour{ШСВ/СВ}
\def\fileExtFpv{lv}
\def\AbzazOne
{
\item
Контроль положения разъединителей при формировании сигналов включенного или отключенного положения выключателей (например, \textit{«В3 включен»} и \textit{«ШСВ/СВ отключен»} для выключателя ШСВ/СВ) вводится программным ключом \textbf{«Контр. Р В3»} (\textbf{«Контр. Р ШСВ/СВ»} для выключателя ШСВ/СВ).
\item
Также функция ФПВ пофазно контролирует состояние блок-контактов выключателя и в случае расхождения полюсов формирует сигнал \textit{«Внутр. пуск. ЗНФ»}.
\item
Функциональная схема логики ФПВ В1 приведена на рисунке \ref{pict_fpv:fpv1}. Уставки ФПВ В1 представлены в таблице \ref{app_set:fpv_q1}.
\item
Функциональная схема логики ФПВ В2 приведена на рисунке \ref{pict_fpv:fpv2}. Описание работы функциональной схемы логики аналогично ФПВ В1. Уставки ФПВ В2 представлены в таблице \ref{app_set:fpv_q2}.
}
}%
{ЮНИТ-М500-ДЗЛ}{%
% код для этого типа
\def\Qthree{В3}
\def\Qfour{ШСВ/СВ}
\def\fileExtFpv{lv}
\def\AbzazOne
{
\item
Контроль положения разъединителей при формировании сигналов включенного или отключенного положения выключателей (например, \textit{«В3 включен»} и \textit{«ШСВ/СВ отключен»} для выключателя ШСВ/СВ) вводится программным ключом \textbf{«Контр. Р В3»} (\textbf{«Контр. Р ШСВ/СВ»} для выключателя ШСВ/СВ).
\item
Также функция ФПВ пофазно контролирует состояние блок-контактов выключателя и в случае расхождения полюсов формирует сигнал \textit{«Внутр. пуск. ЗНФ»}.
\item
Функциональная схема логики ФПВ В1 приведена на рисунке \ref{pict_fpv:fpv1}. Уставки ФПВ В1 представлены в таблице \ref{app_set:fpv_q1}.
\item
Функциональная схема логики ФПВ В2 приведена на рисунке \ref{pict_fpv:fpv2}. Описание работы функциональной схемы логики аналогично ФПВ В1. Уставки ФПВ В2 представлены в таблице \ref{app_set:fpv_q2}.
}
}%
{ЮНИТ-М500-ВЧ}{%
% код для этого типа
\def\Qthree{В3}
\def\Qfour{ШСВ/СВ}
\def\fileExtFpv{lv}
\def\AbzazOne{
\item
Контроль положения разъединителей при формировании сигналов включенного или отключенного положения выключателей (например, \textit{«В3 включен»} и \textit{«ШСВ/СВ отключен»} для выключателя ШСВ/СВ) вводится программным ключом \textbf{«Контр. Р В3»} (\textbf{«Контр. Р ШСВ/СВ»} для выключателя ШСВ/СВ).
\item
Также функция ФПВ пофазно контролирует состояние блок-контактов выключателя и в случае расхождения полюсов формирует сигнал \textit{«Внутр. пуск. ЗНФ»}.
\item
Функциональная схема логики ФПВ В1 приведена на рисунке \ref{pict_fpv:fpv1}. Уставки ФПВ В1 представлены в таблице \ref{app_set:fpv_q1}.
\item
Функциональная схема логики ФПВ В2 приведена на рисунке \ref{pict_fpv:fpv2}. Описание работы функциональной схемы логики аналогично ФПВ В1. Уставки ФПВ В2 представлены в таблице \ref{app_set:fpv_q2}.
}
}%
{ЮНИТ-М500-РЗАТ}{%
% код для этого типа
\def\Qthree{В3}
\def\Qfour{В4}
\def\fileExtFpv{rzat}
\def\AbzazOne{
\item
Контроль положения разъединителей при формировании сигналов включенного или отключенного положения выключателей (например, \textit{«В3 включен»} и \textit{«В3 отключен»} для выключателя В3) вводится программным ключом \textbf{«Контр. Р В3»}. Также имеется возможность выбрать количество разъединителей, участвующих в формировании сигналов включенного или отключенного положения выключателей, для выключателя В3 она задается программным переключателем \textbf{«Схема Р В3»}.
\item
Также функция ФПВ пофазно контролирует состояние блок-контактов выключателя и в случае расхождения полюсов формирует сигнал \textit{«Внутр. пуск. ЗНФ В3»}.
\item
Функциональная схема логики ФПВ В3 приведена на рисунке \ref{pict_fpv:fpv1}. Функциональная схема ФПВ В3 смеж выполняется аналогично. Функциональная схема логики ФПВ В4 приведена на рисунке \ref{pict_fpv:fpv2}. Функциональная схема ФПВ В4 смеж выполняется аналогично. Уставки ФПВ В3 представлены в таблице \ref{app_set:fpv3}. Уставки ФПВ В4 представлены в таблице \ref{app_set:fpv4}. Уставки ФПВ В3 смеж представлены в таблице \ref{app_set:fpv3sm}. Уставки ФПВ В4 смеж представлены в таблице \ref{app_set:fpv4sm}.
}
}%
}
[
\def\Qthree{??}
\def\Qfour{??}
]

\color{unimain}{\subsection{Фиксация положения выключателей  \Qthree~и \Qfour}}
\color{black}

\begin{enumerate}

\item
Функциональный блок фиксации положения выключателей (ФПВ) предназначен для определения состояния выключателей на основании технологической информации:

\begin{itemize}
    \item состояния блок-контактов выключателя;
    \item положения разъединителей в цепи выключателя;
    \item информации о выводе в ремонт выключателя.
\end{itemize}


%%%% СБОРКА СОДЕРЖИМОГО

\AbzazOne

\medskip
\vspace{2mm}
\begin{figure}[H]
\centering
\includegraphics[width=1\textwidth,height=1\textheight,keepaspectratio]{\fbpath/03.01 ЛВ/761.3461 ФПВ34/img/fpv3_\fileExtFpv.pdf}
\captionof{figure}{Функциональная схема логики ФПВ \Qthree}\label{pict_fpv:fpv3}
\end{figure}

\medskip
\vspace{2mm}
\begin{figure}[H]
\centering
\includegraphics[width=1\textwidth,height=1\textheight,keepaspectratio]{\fbpath/03.01 ЛВ/761.3461 ФПВ34/img/fpv4_\fileExtFpv.pdf}
\captionof{figure}{Функциональная схема логики ФПВ \Qfour}\label{pict_fpv:fpv4}
\end{figure}

\end{enumerate}